\documentclass{beamer}

% opacity bugfix: see http://tug.org/pipermail/pdftex/2007-December/007480.html
\pdfpageattr {/Group << /S /Transparency /I true /CS /DeviceRGB>>}

\usepackage[utf8]{inputenc}
\usepackage[OT1]{fontenc}

\usepackage{tikz}
\usetikzlibrary{%
   arrows,%
   calc,%
   fit,%
   patterns,%
   plotmarks,%
   shapes.geometric,%
   shapes.misc,%
   shapes.symbols,%
   shapes.arrows,%
   shapes.callouts,%
   shapes.multipart,%
   shapes.gates.logic.US,%
   shapes.gates.logic.IEC,%
   er,%
   automata,%
   backgrounds,%
   chains,%
   topaths,%
   trees,%
   petri,%
   mindmap,%
   matrix,%
   calendar,%
   folding,%
   fadings,%
   through,%
   patterns,%
   positioning,%
   scopes,%
   decorations.fractals,%
   decorations.shapes,%
   decorations.text,%
   decorations.pathmorphing,%
   decorations.pathreplacing,%
   decorations.footprints,%
   decorations.markings,%
   shadows}
\usepackage{animate}
\usepackage{amssymb, amsmath, amsfonts, enumerate}
\usepackage{pifont}
%\usepackage{bbold}
\newcommand\hmmax{0}
\usepackage{bm}
%\usepackage{dsfont}
\usepackage{pxfonts}
\usepackage{xcolor}
\usepackage{url}

\usepackage[backend=bibtex,style=authoryear,dashed=false]{biblatex}
\addbibresource{refs.bib}
%\renewcommand{\bibfont}{\normalfont\scriptsize}
\setlength{\bibhang}{3ex}

\usepackage{hyperref}

\usetheme[official=false,department=none]{tue2008}
%\usefonttheme{default}


%\usetheme[secheader]{Boadilla}
%\setbeamercovered{transparent}
%\setbeamercovered{invisible}
%\setbeamertemplate{navigation symbols}{}
%\setbeamertemplate{bibliography item}[text] % numbered references
%\useoutertheme{infolines}
%\setbeamertemplate{headline}{}
%\setbeamertemplate{footline}{\hspace*{5mm}\hfill\insertframenumber\hspace*{5mm}\vspace{3mm}}
%\setbeamercolor{alerted text}{fg=orange!80!black}

% ------------------------------------------------------------------------------

\newcommand{\vcenterbox}[1]{\ensuremath{\vcenter{\hbox{#1}}}}

\newcommand{\mbf}[1]{\mathbf{#1}}
\newcommand{\bs}[1]{\boldsymbol{#1}}
\renewcommand{\vec}[1]{{\bm#1}}

\newcommand{\Tsys}{T_\text{sys}}
\newcommand{\Rsys}{R_\text{sys}}
\newcommand{\fsys}{f_\text{sys}}
\newcommand{\Fsys}{F_\text{sys}}

\def\tnow{t_\text{now}}
\def\tpnow{t^+_\text{now}}

\newcommand{\Rsysnow}{R^{(t_\text{now})}_\text{sys}}
\newcommand{\Tsysnow}{T^{(t_\text{now})}_\text{sys}}
\newcommand{\tsysnow}{t^{(t_\text{now})}_\text{sys}}
\newcommand{\fsysnow}{f^{(t_\text{now})}_\text{sys}}
\def\eknow{e_k^{(\tnow)}}
\def\cknow{c_k^{(\tnow)}}
\def\vectknow{\vec{t}_k^{(\tnow)}}
\def\Phinow{\Phi^{(\tnow)}}
\newcommand{\gnow}{g^{(\tnow)}}
\newcommand{\tausnow}{\tau_*^{(\tnow)}}
\newcommand{\tsetup}{\tau_{\text{setup}}}

\newcommand{\cmark}{{\color{tuegreen}\ding{51}}}
\newcommand{\xmark}{{\color{tuered}\ding{55}}}

\newcommand{\comp}[1]{\raisebox{-1mm}{\tikz{\node[%type1
rectangle,rounded corners=0mm,draw,fill=tuepmsgreen!70,thick,inner sep=0pt,minimum size=4mm]{#1};}}}

\newcommand{\cyansec}[1]{\textcolor{tuecyan}{\large\bf #1}}
\newcommand{\cyanalert}[1]{\textcolor{tuecyan}{#1}}
\newcommand{\bluesec}[1]{\textcolor{tueblue}{\large\bf #1}}
\newcommand{\bluealert}[1]{\textcolor{tueblue}{#1}}

\setbeamertemplate{itemize item}{\tiny\raise1.5pt\hbox{\color{tueblue}$\blacktriangleright$}}

% ------------------------------------------------------------------------------


\title{Dynamic and Adaptive Maintenance Policies\\ for Complex Systems}

\author{Gero Walter, Simme Douwe Flapper}
\institute{ Eindhoven University of Technology, Eindhoven, NL\\ 
            \url{g.m.walter@tue.nl} \\[2ex]
            \includegraphics[height=9mm]{logos/tuelogo} \quad 
            \includegraphics[height=9mm]{logos/dinalog-hp} }
\date{CAMPI Slotevent 2016-11-09}

\begin{document}

\frame{
\titlepage
}

\addtocounter{framenumber}{-1}

\begin{frame}{}

\begin{tikzpicture}
[typeA/.style={rectangle,draw,fill=tuepmsgreen!70,very thick,inner sep=0pt,minimum size=6mm},
 typeB/.style={rectangle,draw,fill=tuepmsgreen!70,very thick,inner sep=0pt,minimum size=6mm},
 cross/.style={cross out,draw=red,very thick,minimum width=7mm, minimum height=5mm},
 hv path/.style={very thick, to path={-| (\tikztotarget)}},
 vh path/.style={very thick, to path={|- (\tikztotarget)}}]
\begin{scope}[scale=0.75]
\node[typeA] (A1) at ( 0, 1.2) {A};
\node[typeA] (A2) at ( 0, 0  ) {A};
%\uncover<2->{\node[cross] at (0,0) {};}
\node[typeB] (B1) at ( 0,-1.2) {B};
\coordinate (start) at (-1  ,0);
\coordinate (end)   at ( 1  ,0);
\coordinate (fork1) at (-0.7,0);
\coordinate (fork2) at ( 0.7,0);
\path (fork1) edge[vh path] (start)
              edge[vh path] (A1)
              edge[vh path] (A2)
              edge[vh path] (B1);
\path (fork2) edge[vh path] (end)
              edge[vh path] (A1)
              edge[vh path] (A2)
              edge[vh path] (B1);
\end{scope}
\begin{scope}[xshift=4cm,scale=0.75]
\node[typeA] (A1) at ( 0, 1.2) {A};
%\node[typeA] (A2) at ( 0, 0  ) {A};
%\uncover<2->{\node[cross] at (0,0) {};}
\node[typeB] (B1) at ( 0,-1.2) {B};
\coordinate (start) at (-1  ,0);
\coordinate (end)   at ( 1  ,0);
\coordinate (fork1) at (-0.7,0);
\coordinate (fork2) at ( 0.7,0);
\path (fork1) edge[vh path] (start)
              edge[vh path] (A1)
%             edge[vh path] (A2)
              edge[vh path] (B1);
\path (fork2) edge[vh path] (end)
              edge[vh path] (A1)
%              edge[vh path] (A2)
              edge[vh path] (B1);
\end{scope}
\begin{scope}[yshift=-2cm]
\draw[very thick, -stealth] (-1,0) -- (5,0);
\foreach \t in {0,4} \draw[very thick] (\t,-0.1) -- (\t, 0.1);
\end{scope}
\end{tikzpicture}

\end{frame}

\end{document}
